\subsection{Principal fiber bundles defined by cocycles}

DEFINITION 6. --- Let B be a topological space, let G be a topological group, and let \(\mathcal{U} = (\mathrm{U}_{i})_{i\in \mathrm{I}}\) be an open covering of B. A continuous 1-cocycle on B with values in G, subordinate to \(\mathcal{U}\), is the data, for every pair \((i,j)\) of elements of I, of a continuous map \(g_{i,j}\) from \(\mathrm{U}_i\cap \mathrm{U}_j\) into G, such that for every triple \((i,j,k)\in \mathrm{I}\times \mathrm{I}\times \mathrm{I}\) and every point b of \(\mathrm{U}_i\cap \mathrm{U}_j\cap \mathrm{U}_k\), one has

\[
g _ {i, k} (b) = g _ {i, j} (b) g _ {j, k} (b).
\]

We denote \(Z_{\mathrm{cont}}^{1}(\mathrm{B},\mathcal{U},\mathrm{G})\) the set of continuous 1-cocycles on B with values in G, subordinate to the covering U.

In this issue, we will simply say cocycle instead of continuous 1-cocycle.

Let \((\mathrm{E},p)\) be a principal G-bundle over B and let \(\mathcal{U} = (\mathrm{U}_{i})_{i\in\mathrm{I}}\) be a covering of B by open sets \(U_{i}\). We say that E is trivializable over U if, for every \(i\in I\), E is trivializable over \(U_{i}\). We then call a U-trivialization of E a family \((f_{i})_{i\in\mathrm{I}}\), where \(f_{i}\colon p^{-1}(\mathrm{U}_{i})\to\mathrm{U}_{i}\times\mathrm{G}\) is a trivialization of E over \(U_{i}\).

Let \(\mathcal{U} = (\mathrm{U}_i)_{i\in \mathrm{I}}\) be a covering of B by open sets, and let \((\mathrm{E},p)\) be a principal G-bundle over B equipped with a \(\mathcal{U}\)-trivialization \((f_i)_{i\in \mathrm{I}}\). For every pair \((i,j)\in \mathrm{I}\times \mathrm{I}\), denote by \(g_{ij}\) the map from \(\mathrm{U}_i\cap \mathrm{U}_j\) into G defined by

\[
(b, g _ {i j} (b)) = f _ {i} \circ f _ {j} ^ {- 1} (b, e),
\]

where e denotes the identity element of G. It is continuous.

Since \(f_{i}\) and \(f_{j}\) are compatible with the operations of G, we have

\[
(f _ {i} \circ f _ {j} ^ {- 1}) (b, g) = (b, g _ {i j} (b) g)
\]

for \(b \in \mathrm{U}_i \cap \mathrm{U}_j\) and \(g \in \mathrm{G}\). If \(b\) is a point of \(\mathrm{U}_i \cap \mathrm{U}_j \cap \mathrm{U}_k\) (where \(i, j, k \in \mathrm{I}\)), we have

\[
(f _ {i} \circ f _ {k} ^ {- 1}) (b, e) = (b, g _ {i k} (b))
\]

and

\[
\begin{array}{c} (f _ {i} \circ f _ {k} ^ {- 1}) (b, e) = (f _ {i} \circ f _ {j} ^ {- 1}) \circ (f _ {j} \circ f _ {k} ^ {- 1}) (b, e) \\ = (f _ {i} \circ f _ {j} ^ {- 1}) (b, g _ {j k} (b)) = (b, g _ {i j} (b) g _ {j k} (b)). \end{array}
\]

It follows that

\[
g _ {i k} (b) = g _ {i j} (b) g _ {j k} (b),
\]

so that the family \((g_{ij})\) is a cocycle on B with values in G, subordinate to the covering U. It is called the cocycle defined by the family of trivializations \((f_{i})_{i\in\mathrm{I}}\).

Let B, G, and U be as in Definition 6, and let \((g_{ij}) \in Z_{\mathrm{cont}}^{1}(\mathrm{B}, \mathcal{U}, \mathrm{G})\) be a cocycle. For every pair \((i, j) \in \mathrm{I} \times \mathrm{I}\), the map \(\gamma_{ij} \colon (\mathrm{U}_{i} \cap \mathrm{U}_{j}) \times \mathrm{G} \to (\mathrm{U}_{i} \cap \mathrm{U}_{j}) \times \mathrm{G}\) defined by

\[
\gamma_ {i j} (b, g) = \left(b, g _ {i j} (b) g\right) \quad \text {   for   } b \in \mathrm{U} _ {i} \cap \mathrm{U} _ {j} \text {   and   } g \in \mathrm{G}\tag{5}
\]

is an isomorphism of principal bundles over the base \(U_{i} \cap U_{j}\) . For every triple \((i,j,k) \in \mathrm{I} \times \mathrm{I} \times \mathrm{I}\) and every pair \((b,g) \in (\mathrm{U}_{i} \cap \mathrm{U}_{j} \cap \mathrm{U}_{k}) \times \mathrm{G}\) , we have:

\[
\gamma_ {i k} (b, g) = \gamma_ {i j} \circ \gamma_ {j k} (b, g).
\]

Let F be the topological space obtained by gluing together the spaces \(U_{i} \times G\) along the \((\mathrm{U}_{i} \cap \mathrm{U}_{j}) \times \mathrm{G}\) by means of the bijections \(\gamma_{ij}\) (TG, I, p. 16). For every \(i \in I\), the image of \(U_{i} \times G\) in F is an open set of F (TG, I, p. 17, prop. 9). The canonical projections \(p_{i}: U_{i} \times G \to U_{i}\) glue together into a continuous map p from F to B. The maps \(\gamma_{ij}\) being compatible with the right G-operations in the spaces \(U_{i} \times G\), one deduces a continuous right action of G on F which makes F a principal fiber bundle with base B and group G, equipped with a trivialization over each \(U_{i}\). We say that F is the principal fiber bundle defined by the cocycle \((g_{ij})\).

DEFINITION 7. --- Let B be a topological space, let G be a topological group, and let \(\mathcal{U} = (\mathrm{U}_{i})_{i\in \mathrm{I}}\) be an open covering of B. We say that two cocycles \((g_{ij})\) and \((g_{ij}^{\prime})\) in \(\mathrm{Z}_{\mathrm{cont}}^{1}(\mathrm{B},\mathcal{U},\mathrm{G})\) are cohomologous if there exists a family \((h_{i})_{i\in\mathrm{I}}\) of continuous maps \(h_{i}\colon U_{i}\to G\) such that one has

\[
g _ {i j} ^ {\prime} (b) = h _ {i} (b) g _ {i j} (b) h _ {j} (b) ^ {- 1}\tag{6}
\]

for every pair \((i,j)\in \mathrm{I}\times \mathrm{I}\) and every \(b\in \mathrm{U}_i\cap \mathrm{U}_j\).

The relation “\((g_{ij})\) is cohomologous to \((g'_{ij})\)” is an equivalence relation on the set \(Z_{\mathrm{cont}}^{1}(\mathrm{B},\mathcal{U},\mathrm{G})\). We denote by \(H_{\mathrm{cont}}^{1}(\mathrm{B},\mathcal{U},\mathrm{G})\) the quotient set of \(Z_{\mathrm{cont}}^{1}(\mathrm{B},\mathcal{U},\mathrm{G})\) by this equivalence relation.

PROPOSITION 11. --- Let B be a topological space, G a topological group, and \(\mathcal{U} = (\mathrm{U}_{i})_{i\in \mathrm{I}}\) an open covering of B.

\begin{enumerate}
\def\labelenumi{\alph{enumi})}
\item
  Every principal G-bundle over B that is trivializable over U is isomorphic to a principal bundle defined by a cocycle in \(Z_{\mathrm{cont}}^{1}(\mathrm{B}, \mathcal{U}, \mathrm{G})\).
\item
  Let (E, p) and (E', p') be principal G-bundles over B that are trivializable over U. Let \((f_i)_{i \in I}\) (resp. \((f_i')_{i \in I}\)) be a trivialization of (E, p) (resp. of \((E', p')\)) adapted to U, and denote by \((g_{ij})_{(i,j) \in I \times I}\) (resp. \((g_{i,j}')_{(i,j) \in I \times I}\)) the cocycle defined by this trivialization. Then, the principal fiber bundles (E, p) and (E', p') are isomorphic if and only if these cocycles are cohomologous.
\end{enumerate}

Let us prove \(b\). Let \(\varphi: \mathrm{E} \to \mathrm{E}'\) be an isomorphism of principal bundles with base B and group G. For every \(i \in \mathrm{I}\), let \(h_i\) be the continuous map from \(\mathrm{U}_i\) into G defined by

\[
(b, h _ {i} (b)) = f _ {i} ^ {\prime} \circ \varphi \circ f _ {i} ^ {- 1} (b, e).
\]

Since, for every \(i \in I\) , \(f_{i}\) and \(f_{i}^{\prime}\) are compatible with the operations of G, we have for every \(b \in U_{i}\) and every \(g \in G\) ,

\[
(f _ {i} ^ {\prime} \circ \varphi \circ f _ {i} ^ {- 1}) (b, g) = (b, h _ {i} (b) g)
\]

and

\[
(f _ {i} \circ \varphi^ {- 1} \circ (f _ {i} ^ {\prime}) ^ {- 1}) (b, g) = (b, h _ {i} (b) ^ {- 1} g).
\]

Consequently, for every pair \((i,j)\in \mathrm{I}\times \mathrm{I}\) and every point \(b\) of \(\mathrm{U}_i\cap \mathrm{U}_j\),

\[
\begin{array}{c} f _ {i} ^ {\prime} \circ (f _ {j} ^ {\prime}) ^ {- 1} (b, e) = (f _ {i} ^ {\prime} \circ \varphi \circ f _ {i} ^ {- 1}) \circ (f _ {i} \circ f _ {j} ^ {- 1}) \circ (f _ {j} \circ \varphi^ {- 1} \circ (f _ {j} ^ {\prime}) ^ {- 1}) (b, e) \\ = (b, h _ {i} (b) g _ {i j} (b) h _ {j} (b) ^ {- 1}) \end{array}
\]

so that we have

\[
g _ {i j} ^ {\prime} (b) = h _ {i} (b) g _ {i j} (b) h _ {j} (b) ^ {- 1}.
\]

This proves that the cocycles \((g_{ij})\) and \((g'_{ij})\) are cohomologous.

Conversely, suppose that these cocycles are cohomologous, and let \((h_{i})_{i\in\mathrm{I}}\) be a family of continuous maps \(h_{i}\colon U_{i}\to G\) such that one has \(g_{ij}^{\prime}(b)=h_{i}(b)g_{ij}(b)h_{j}(b)^{-1}\) for \(i,j\in I\) and \(b\in U_{i}\cap U_{j}\). For \(i\in I\), let \(\varphi_{i}\colon p^{-1}(U_{i})\to(p')^{-1}(U_{i})\) be the map defined by \(f_{i}^{\prime}\circ\varphi_{i}\circ f_{i}^{-1}(b,g)=(b,h_{i}(b)g)\), for \(b\in U_{i}\) and \(g\in G\). This is an isomorphism of principal fiber bundles with base \(U_{i}\) and group G. For \((i,j)\in I\times I\), denote by \(\gamma_{ij}\) and \(\gamma_{ij}^{\prime}\) the gluing homeomorphisms associated as above to the cocycles \((g_{ij})\) and \((g_{ij}^{\prime})\), respectively (I, p. 115, formula (5)), so that \(f_{i}(b,g)=\gamma_{ij}\circ f_{j}(b,g)\) and \(f_{i}^{\prime}(b,g)=\gamma_{ij}^{\prime}\circ f_{j}^{\prime}(b,g)\) for all \((b,g)\in(U_{i}\cap U_{j})\times G\). Consequently, for \((i,j)\in I\times I\) and \((b,g)\in(U_{i}\cap U_{j})\times G\), we have the relations

\[
\begin{array}{r l} & f _ {i} ^ {\prime} \circ \varphi_ {j} \circ f _ {i} ^ {- 1} (b, g) = \gamma_ {i j} ^ {\prime} \circ (f _ {j} ^ {\prime} \circ \varphi_ {j} \circ f _ {j} ^ {- 1}) (b, g _ {i j} (b) ^ {- 1} g) \\ & \qquad = \gamma_ {i j} ^ {\prime} (b, h _ {j} (b) g _ {i j} (b) ^ {- 1} g) \\ & \qquad = (b, g _ {i j} ^ {\prime} (b) h _ {j} (b) g _ {i j} (b) ^ {- 1} g) \\ & \qquad = (b, h _ {i} (b) g) = f _ {i} ^ {\prime} \circ \varphi_ {i} \circ f _ {i} ^ {- 1} (b, g). \end{array}
\]

This demonstrates that \(\varphi_{i}\) and \(\varphi_{j}\) coincide on \(p^{-1}(\mathrm{U}_{i}\cap\mathrm{U}_{j})\). The morphisms \(\varphi_{i}\) therefore glue together into a B-morphism of principal bundles from E to E'. The assertion \(b\) follows because any morphism of principal bundles with base B and group G is an isomorphism (I, p. 93, prop. 1).

Now let us prove \(a\). Let \((\mathrm{E}, p)\) be a principal fiber bundle with structure group G, equipped, for every \(i \in \mathrm{I}\), with a trivialization \(f_i: p^{-1}(\mathrm{U}_i) \to \mathrm{U}_i \times \mathrm{G}\) over \(\mathrm{U}_i\). Let \((g_{ij})\) be the cocycle defined by this family, and let F be the principal fiber bundle defined by the cocycle \((g_{ij})\). By construction, the principal fiber bundle F is equipped with a trivialization over \(\mathcal{U}\); this trivialization defines the cocycle \((g_{ij})\). By assertion \(b\), the principal fiber bundle \((\mathrm{E}, p)\) is isomorphic to F.

Let B be a topological space and G a topological group. Let \((E,p)\) be a principal fiber bundle with base B and group G. Let s be a section of the surjective map p (E, II, p. 19, prop. 8). The map \(f\colon B\times G\to E\) defined by \(f(b,g)=s(b)\cdot g\) is a bijection compatible with the action of G. Equip \(B\times G\) with the topology obtained by transport of structure; \((B\times G,\mathrm{pr}_{1})\) is then a principal fiber bundle with base B and group G isomorphic to E. There therefore exists a set T of principal bundles with base B and group G such that every principal fiber bundle with base B and group G is isomorphic to an element of T. Let P(B, G) denote the set of isomorphism classes of principal bundles with base B and group G (E, II, p. 47).

Let \(\mathcal{U} = (\mathrm{U}_i)_{i\in \mathrm{I}}\) be an open cover of B. Let us denote by \(\mathrm{P(B,\mathcal{U},G)}\) the subset of \(\mathrm{P(B,G)}\) consisting of the isomorphism classes of principal bundles that are trivializable over \(\mathcal{U}\). By Proposition 11, there exists a map \(r_{\mathcal{U}}\) from \(\mathrm{H}_{\mathrm{cont}}^{1}(\mathrm{B},\mathcal{U},\mathrm{G})\) into \(\mathrm{P(B,\mathcal{U},G)}\) which associates with the class of a cocycle \((g_{ij})\in Z_{\mathrm{cont}}^{1}(\mathrm{B},\mathcal{U},\mathrm{G})\) the isomorphism class of the principal bundle defined by this cocycle; it is a bijection (loc. cit.). The inverse bijection \(s_{\mathcal{U}}\) associates with the isomorphism class in \(\mathrm{P(B,\mathcal{U},G)}\) of a principal fiber bundle E, trivializable over \(\mathcal{U}\), the class of the cocycle defined by any family of trivializations of E over \(\mathcal{U}\). In this correspondence, the isomorphism class of the trivial principal fiber bundle \(\mathrm{B}\times \mathrm{G}\) corresponds to the cohomology class of the constant cocycle \(g_{ij} = e\), also called the trivial cocycle.

By virtue of the definition of a principal fiber bundle, the set P(B, G) is the union of the sets of the form P(B, U, G), where U is an open covering of B.

Let \(\mathcal{U} = (\mathrm{U}_i)_{i\in \mathrm{I}}\) be an open cover of B and let \(\mathcal{V} = (\mathrm{V}_k)_{k\in \mathrm{K}}\) be an open cover of B finer than \(\mathcal{U}\). We have \(\mathrm{P(B,\mathcal{U},G)}\subset\) \(\mathrm{P(B,\mathcal{V},G)}\); we denote by \(i_{\mathcal{V}\mathcal{U}}\) the canonical injection defined by this inclusion. Let us choose a map \(\varphi \colon \mathrm{K}\to \mathrm{I}\) such that \(\mathrm{V}_k\subset \mathrm{U}_{\varphi (k)}\) for all \(k\in \mathrm{K}\). Given a cocycle \((g_{ij})\in \mathrm{Z}_{\mathrm{cont}}^1 (\mathrm{B},\mathcal{U},\mathrm{G})\), set, for every pair \((k,\ell)\in \mathrm{K}\times \mathrm{K}\), \(\overline{g}_{k\ell} = g_{\varphi (k)\varphi (\ell)}\mid \mathrm{V}_k\cap \mathrm{V}_\ell\).

The family \((\overline{g}_{k\ell})\) is a cocycle on B, with values in G, subordinate to V. If \((g'_{ij}) \in \mathrm{Z}_{\mathrm{cont}}^{1}(\mathrm{B}, \mathcal{U}, \mathrm{G})\) is a cocycle cohomologous to the cocycle \((g_{ij})\), the cocycle \((\overline{g}'_{k\ell})\) deduced from \((g'_{ij})\) is cohomologous to the cocycle \((\overline{g}_{k\ell})\). This yields a map

\[
c (\varphi) \colon \mathrm{H} _ {\mathrm{cont}} ^ {1} (\mathrm{B}, \mathcal {U}, \mathrm{G}) \to \mathrm{H} _ {\mathrm{cont}} ^ {1} (\mathrm{B}, \mathcal {V}, \mathrm{G})
\]

which to the class of \((g_{ij})\) associates the class of \((\overline{g}_{k\ell})\) .

Let E be a principal fiber bundle with base B and group G and, for every \(i \in I\), let \(f_i: E_{U_i} \to U_i \times G\) be a trivialization of the \(U_i\)-principal fiber bundle \(E_{U_i}\). For all \(k \in K\), let \(f_k' : E_{V_k} \to V_k \times G\) be the trivialization deduced from \(f_{\varphi(k)}\) by passing to subsets. Let \((g_{ij}) \in Z_{\mathrm{cont}}^1(\mathrm{B}, \mathcal{U}, \mathrm{G})\) be the cocycle defined by the family \((f_i)\); the cocycle defined by the family \((f_{k}^{\prime})\) is precisely the cocycle \((\overline{g}_{k\ell})\) defined above. Thus, the following diagram is commutative:

\[
\begin{array}{c} \mathrm{H} _ {\mathrm{cont}} ^ {1} (\mathrm{B}, \mathcal {U}, \mathrm{G}) ^ {c (\varphi)} \longrightarrow \mathrm{H} _ {\mathrm{cont}} ^ {1} (\mathrm{B}, \mathcal {V}, \mathrm{G}) \\ \Big \downarrow^ {r _ {\mathcal {U}}} \qquad \qquad \qquad \qquad \Big \downarrow^ {r _ {\mathcal {V}}} \\ \mathrm{P} (\mathrm{B}, \mathcal {U}, \mathrm{G}) ^ {i _ {\mathcal {V} \mathcal {U}}} \longrightarrow \mathrm{P} (\mathrm{B}, \mathcal {V}, \mathrm{G}). \end{array}
\]

The maps \(r_{\mathcal{U}}\) and \(r_{\mathcal{V}}\) being bijective, the map \(c(\varphi)\), just as \(i_{\mathcal{V}\mathcal{U}}\), is an injection and does not depend on the choice of \(\varphi\). Henceforth we shall write \(c_{\mathcal{V}\mathcal{U}}\) instead of \(c(\varphi)\).

Let \(\mathcal{R}\) be the set of elements of \(\mathfrak{P}(\mathfrak{P}(B))\) that are open coverings of B. For every open covering \(\mathcal{U}\) of B, there exists an open covering \(\mathcal{V}\) belonging to \(\mathcal{R}\) such that \(\mathcal{V}\) is both finer and coarser than \(\mathcal{U}\). The set \(\mathcal{R}\) is ordered and filtering for the relation \(\leqslant\) defined by \(\mathcal{U} \leqslant \mathcal{V}\) if \(\mathcal{V}\) is a finer covering than \(\mathcal{U}\). It follows from the preceding that one has defined an inductive system \((\mathrm{H}_{\mathrm{cont}}^{1}(\mathrm{B},\mathcal{U},\mathrm{G}),c_{\mathcal{V}\mathcal{U}})\) relative to the filtered ordered set \(\mathcal{R}\) and that the family \((r_{\mathcal{U}})\) is an inductive system of bijective maps from \((\mathrm{H}_{\mathrm{cont}}^{1}(\mathrm{B},\mathcal{U},\mathrm{G}),c_{\mathcal{V}\mathcal{U}})\) into \((\mathrm{P}(\mathrm{B},\mathcal{U},\mathrm{G}),i_{\mathcal{V}\mathcal{U}})\). If we denote by \(\mathrm{H}_{\mathrm{cont}}^{1}(\mathrm{B},\mathrm{G})\) the inductive limit of the system \((\mathrm{H}_{\mathrm{cont}}^{1}(\mathrm{B},\mathcal{U},\mathrm{G}),c_{\mathcal{V}\mathcal{U}})\) and \(r\colon \mathrm{H}_{\mathrm{cont}}^{1}(\mathrm{B},\mathrm{G})\to \mathrm{P}(\mathrm{B},\mathrm{G})\) the inductive limit of the family \((r_{\mathcal{U}})\), we therefore have:

THEOREM 4. --- The map \(r\colon \mathrm{H}_{\mathrm{cont}}^{1}(\mathrm{B},\mathrm{G})\to \mathrm{P}(\mathrm{B},\mathrm{G})\) is bijective.

Let \(\mathcal{U} = (\mathrm{U}_i)_{i\in \mathrm{I}}\) be an open covering of B; let us denote by \(c_{\mathcal{U}}\) the canonical map \(\mathrm{H}_{\mathrm{cont}}^{1}(\mathrm{B},\mathcal{U},\mathrm{G})\to \mathrm{H}_{\mathrm{cont}}^{1}(\mathrm{B},\mathrm{G})\). If \((g_{ij})\) is a cocycle on B, with values in G, subordinate to the open covering \(\mathcal{U}\), the element \(c_{\mathcal{U}}((g_{ij}))\) of \(\mathrm{H}_{\mathrm{cont}}^{1}(\mathrm{B},\mathrm{G})\) is called the cohomology class of the cocycle \((g_{ij})\).

